\chapter*{Declaração de uso de assistência de inteligência artificial}
\addcontentsline{toc}{chapter}{Declaração de uso de assistência de IA}

Este trabalho fez uso de assistência de inteligência artificial
(Claude, Anthropic) como ferramenta de apoio, sob direção e revisão
contínuas do autor, nas seguintes etapas: redação assistida do texto a
partir de resultados, esboços e decisões do autor; implementação
assistida de código de experimentos e de análise; verificações
cruzadas, repetidas e datadas, entre números do texto e os artefatos
versionados que os sustentam; apoio em buscas bibliográficas e nas
verificações datadas de novidade (com conferência das fontes
primárias); e organização de material de apresentação. O desenho de
pesquisa, as hipóteses, as decisões metodológicas, a seleção do que
reportar --- incluindo os resultados negativos --- e a validação final
de cada afirmação são do autor.

Além do apoio à escrita, o texto foi submetido, em 23/11/2026, a uma
\textbf{revisão adversarial automatizada}: uma sessão independente de
IA, sem acesso ao contexto de escrita, instruída a atuar como banca
examinadora --- recomputando números a partir dos artefatos brutos
(cerca de 400 verificações), reexecutando experimentos e confrontando
texto, código e registro público. As correções decorrentes dessa
revisão estão incorporadas e registradas no adendo de 23/11/2026 ao
documento de desenho experimental e no histórico do repositório; as
verificações centrais foram re-executadas a partir dos artefatos
brutos, nesta mesma infraestrutura assistida, antes de serem
incorporadas.

Nenhum resultado experimental foi gerado por modelo de linguagem:
todos os números provêm de artefatos produzidos pelos scripts
versionados do repositório público do projeto, com sementes fixadas e
execução reproduzível. O histórico de desenvolvimento é público no
repositório, parte dos \emph{commits} registra explicitamente a
coautoria da ferramenta, e as correções de registro foram feitas por
retificação datada, nunca por reescrita de histórico.
